%===================================================================
\appendix
\section{Digital-Twin Construction for Ferritic--Martensitic Steel Microstructure}
\label{app:radcluster_digital_twin}
%===================================================================

This appendix gives a detailed plan for constructing a RadCluster-based
\emph{digital twin} of irradiated ferritic--martensitic (F/M) steel
microstructure.  The twin is intended to transform the existing experimental
microstructure database into a calibrated, uncertainty-aware predictive model
for loop and cavity evolution under fission, ion, and ultimately fusion-relevant
irradiation conditions.  The guiding principle is that the digital twin is not a
single optimized RadCluster calculation; it is a probability distribution over
RadCluster inputs and predictions, constrained by measured loop and cavity
statistics and by physical conservation laws.

Digital twins have increasingly been formulated as hybrid physics--data systems
whose predictive value depends on uncertainty quantification (UQ), calibration,
validation, and sequential updating.  Recent reviews emphasize that uncertainty
quantification and optimization are enabling technologies for digital-twin
credibility, especially when the physical twin is only partially observable and
when the simulator must extrapolate beyond the measured data
\citep{Thelen2022DigitalTwinUQ}.  In materials science, related work has used
Bayesian learning and physics-informed surrogate models to combine simulations
with sparse experiments, calibrate uncertain parameters, and propagate
prediction uncertainty in additive-manufacturing and microstructure-sensitive
applications \citep{Wang2020AMUQ,Kim2022BayesianAM,Robertson2026VDMN}.
The present RadCluster twin follows the same philosophy but applies it to the
irradiation microstructure of F/M steels: the forward model is the cluster
dynamics simulator, the uncertain parameters are physically bounded production,
transport, binding, sink, and loop-conversion quantities, and the data are the
measured loop and cavity densities, mean diameters, and size distributions.

The need for this construction follows directly from the main text.  RadCluster
is already formulated as a size- and composition-resolved cluster dynamics model
for SIA clusters, vacancy clusters/cavities, helium--vacancy complexes, loop
character conversion, sinks, and conservation diagnostics.  Section~2.4 defines
a UQ workflow based on space-filling design, Sobol sensitivity screening,
multi-fidelity Gaussian-process emulation, Bayesian calibration, Wasserstein
comparison of size distributions, and cost-aware active refinement.  Appendix~G
then supplies the statistical target: an experimental microstructure database
containing neutron and ion irradiation data for loop number density, mean loop
diameter, cavity number density, mean cavity diameter, and selected loop-size
distributions.  The purpose of this appendix is to turn those ingredients into a
concrete digital-twin algorithm.

%===================================================================
\subsection{Definition of the Digital Twin}
\label{app:dt_definition}
%===================================================================

Let the RadCluster state at dose $d$ and temperature $T$ be
\begin{equation}
    \mathbf{x}(d,T;\boldsymbol{\theta})
    =
    \left[
      \{c^I_{n,\alpha}\},
      \{c^V_{m,\boldsymbol{\ell},\beta}\},
      \{c^g_\mu\},
      \rho_{\mathrm{net}},
      \Delta J_i^{\mathrm{fix}},
      \Delta J_v^{\mathrm{fix}},
      \Delta J_g^{\mathrm{fix}}
    \right]^{\top},
    \label{eq:dt_state}
\end{equation}
where $c^I_{n,\alpha}$ are interstitial-cluster populations,
$c^V_{m,\boldsymbol{\ell},\beta}$ are vacancy-cluster/cavity populations,
$c^g_\mu$ are matrix gas reservoirs, $\rho_{\mathrm{net}}$ is the evolving
dislocation-network density when activated, and the $\Delta J$ variables are
cumulative sink-flux diagnostics.

The uncertain RadCluster parameter vector is denoted by
\begin{equation}
    \boldsymbol{\theta}
    =
    \left(
      \theta_1,\theta_2,\ldots,\theta_p
    \right),
    \qquad p\simeq 16\text{--}20 .
    \label{eq:dt_theta}
\end{equation}
For a given irradiation condition
\begin{equation}
    \boldsymbol{\xi}
    =
    \left(T,d,\dot d,\eta_{\mathrm{He}},\eta_{\mathrm{H}},
    \mathrm{alloy},\mathrm{irradiation\ type}\right),
\end{equation}
where $\eta_{\mathrm{He}}$ and $\eta_{\mathrm{H}}$ are the gas-production ratios
in appm/dpa, RadCluster defines a deterministic map
\begin{equation}
    \mathbf{y}_{\mathrm{sim}}
    =
    H\!ig[\mathbf{x}(d,T;\boldsymbol{\theta})\big]
    =
    \mathcal{M}(\boldsymbol{\theta};\boldsymbol{\xi}),
    \label{eq:dt_forward_map}
\end{equation}
where $H$ is the observation operator that maps the full cluster population
to experimentally measured quantities.  The standard observable vector is
\begin{equation}
    \mathbf{y}_{\mathrm{sim}}
    =
    \left[
      N_\ell,
      \bar d_\ell,
      N_c,
      \bar d_c,
      f_{111},
      f_{100},
      S,
      p_\ell(d_\ell),
      p_c(d_c)
    \right].
    \label{eq:dt_observables}
\end{equation}
Here $N_\ell$ and $\bar d_\ell$ are the loop number density and mean loop
diameter, $N_c$ and $\bar d_c$ are the cavity number density and mean cavity
diameter, $f_{111}$ and $f_{100}$ are loop-character fractions, $S$ is swelling,
and $p_\ell$ and $p_c$ are loop and cavity size distributions.

The digital twin is defined by the posterior distribution
\begin{equation}
    p(\boldsymbol{\theta}\mid \mathcal{D}_{\mathrm{exp}})
    \propto
    p(\mathcal{D}_{\mathrm{exp}}\mid \boldsymbol{\theta})
    p(\boldsymbol{\theta}),
    \label{eq:dt_posterior}
\end{equation}
and by the posterior predictive distribution
\begin{equation}
    p(\mathbf{y}_{\ast}\mid\mathcal{D}_{\mathrm{exp}},\boldsymbol{\xi}_\ast)
    =
    \int
    p(\mathbf{y}_{\ast}\mid\boldsymbol{\theta},\boldsymbol{\xi}_\ast)
    p(\boldsymbol{\theta}\mid\mathcal{D}_{\mathrm{exp}})
    \,d\boldsymbol{\theta} .
    \label{eq:dt_posterior_predictive}
\end{equation}
Thus, the twin reports medians, credible intervals, and risk probabilities,
not only best-fit predictions.

%===================================================================
\subsection{Uncertain Parameter Vector}
\label{app:dt_parameter_vector}
%===================================================================

The first-generation twin should use a restricted but physically interpretable
parameter vector.  A representative 18-parameter set is
\begin{equation}
\begin{split}
\boldsymbol{\theta}=ig[&
\eta_{\mathrm{FP}},
 f_I^{\mathrm{cl}},
 f_V^{\mathrm{cl}},
 E_m^I,
 E_m^V,
 E_m^{\mathrm{He}},
 L_{1D},
 i_{\mathrm{mob}},
 v_{\mathrm{mob}},
 A_V,
 B_V,
 A_I,
 B_I,
 A_{\mathrm{HeV}},
 Z_i^d,
 Z_i^{\mathrm{loop}},
 T^*,
 K_{\mathrm{rec}}
\big].
\end{split}
\label{eq:dt_theta_18}
\end{equation}
The parameters fall into five mechanism groups:
\begin{enumerate}
    \item \emph{Production:} cascade survival efficiency and fractions of
    SIAs and vacancies born as clusters.
    \item \emph{Mobility:} migration energies, the effective one-dimensional
    glide length, and the maximum mobile cluster sizes.
    \item \emph{Dissociation energetics:} binding-law amplitudes and exponents
    controlling SIA, vacancy, and gas emission.
    \item \emph{Sinks and bias:} dislocation bias and loop bias factors.
    \item \emph{Loop conversion and network evolution:} the
    $\frac{1}{2}\langle111\rangle\rightarrow\langle100\rangle$ crossover
    temperature and the dislocation-network recovery constant.
\end{enumerate}

Each parameter is assigned a bounded prior distribution.  Energies are most
naturally represented by truncated normal priors,
\begin{equation}
    E_i \sim \mathcal{N}_{[E_i^{-},E_i^{+}]}
    \left(\bar E_i,\sigma_{E_i}^2\right),
\end{equation}
positive rate or density parameters by log-normal priors,
\begin{equation}
    \rho_d,
    K_{\mathrm{rec}},
    L_{1D}
    \sim
    \mathrm{Lognormal}(\mu,\sigma^2),
\end{equation}
and bounded fractions by beta or uniform priors,
\begin{equation}
    f_I^{\mathrm{cl}}, f_V^{\mathrm{cl}}
    \sim
    \mathrm{Beta}(a,b)
    \quad\mathrm{or}\quad
    \mathcal{U}(0,1).
\end{equation}
Discrete quantities such as $i_{\mathrm{mob}}$, $v_{\mathrm{mob}}$, and possible
conversion-onset sizes should be treated either as categorical model variants
or as integer-valued sampled parameters.

%===================================================================
\subsection{Experimental Target Representation}
\label{app:dt_targets}
%===================================================================

Let the experimental database be
\begin{equation}
    \mathcal{D}_{\mathrm{exp}}
    =
    \left\{
      \boldsymbol{\xi}_j,
      \mathbf{z}_j,
      \boldsymbol{\sigma}_j,
      g_j
    \right\}_{j=1}^{N_{\mathrm{data}}},
\end{equation}
where $\boldsymbol{\xi}_j$ contains irradiation conditions, $\mathbf{z}_j$
contains measured observables, $\boldsymbol{\sigma}_j$ contains measurement
or fitted statistical uncertainties, and $g_j$ is the grouping variable used in
Appendix~G, such as irradiation type, alloy family, dose group, or
bubble/void classification.

For scalar positive observables, Appendix~G uses log-linear statistical fits,
\begin{equation}
    \ln y
    =
    a
    +
    \sum_g c_g\delta_g
    +
    b x
    +
    \epsilon,
    \qquad
    \epsilon\sim\mathcal{N}(0,s^2),
    \label{eq:dt_database_fit}
\end{equation}
where $x$ is either $T_{\mathrm{eq}}$ or $1000/T_{\mathrm{eq}}$.  Therefore, the
natural comparison between RadCluster and the database is in log space:
\begin{equation}
    r_j(\boldsymbol{\theta})
    =
    \frac{
    \ln y^{\mathrm{sim}}_j(\boldsymbol{\theta})
    -
    \ln y^{\mathrm{exp}}_j
    }{
    \sigma_{\ln y,j}
    } .
    \label{eq:dt_log_residual}
\end{equation}
When the comparison is made to the fitted trend rather than to an individual
experimental row,
\begin{equation}
    y_{j,\mathrm{fit}}
    =
    \exp(a+c_{g,j}+b x_j),
    \qquad
    \sigma_{\ln y,j}=s_j .
\end{equation}
A simulation is statistically consistent with the fitted database band if
\begin{equation}
    \left|\ln y^{\mathrm{sim}}_j-\ln y_{j,\mathrm{fit}}\right|
    \lesssim
    2s_j .
    \label{eq:dt_band_acceptance}
\end{equation}
Equivalently, $y_j^{\mathrm{sim}}$ lies within the multiplicative interval
$[e^{-2s_j},e^{2s_j}]$ about the fitted median.

For full loop or cavity size distributions, the discrepancy should not be
reduced to mean diameter alone.  Let $p_j^{\mathrm{exp}}(R)$ and
$p_j^{\mathrm{sim}}(R;\boldsymbol{\theta})$ be the measured and simulated size
distributions.  The distributional contribution to the likelihood is measured
by an optimal-transport distance,
\begin{equation}
    \mathcal{W}_j(\boldsymbol{\theta})
    =
    W_1\!ig[
      p_j^{\mathrm{sim}}(R;\boldsymbol{\theta}),
      p_j^{\mathrm{exp}}(R)
    \big],
    \label{eq:dt_wasserstein}
\end{equation}
which penalizes both shifts in the location of the distribution and changes in
its width or shape.

%===================================================================
\subsection{Likelihood, Model Discrepancy, and Conservation Penalties}
\label{app:dt_likelihood}
%===================================================================

The experimental observation model is written as
\begin{equation}
    \ln y_j^{\mathrm{exp}}
    =
    \ln y_j^{\mathrm{sim}}(\boldsymbol{\theta})
    +
    \delta_j^{\mathrm{model}}
    +
    \epsilon_j^{\mathrm{exp}},
    \label{eq:dt_obs_model}
\end{equation}
where $\epsilon_j^{\mathrm{exp}}$ is measurement noise and
$\delta_j^{\mathrm{model}}$ is structural model discrepancy.  The total variance
used in calibration is
\begin{equation}
    \sigma_j^2
    =
    \sigma_{\mathrm{exp},j}^2
    +
    \sigma_{\mathrm{db},j}^2
    +
    \sigma_{\mathrm{em},j}^2
    +
    \sigma_{\mathrm{model},j}^2,
    \label{eq:dt_uncertainty_sum}
\end{equation}
where the four terms represent measurement uncertainty, database scatter,
emulator uncertainty, and irreducible model-form discrepancy.

The negative log-likelihood is then
\begin{equation}
\begin{split}
    -\log p(\mathcal{D}_{\mathrm{exp}}\mid\boldsymbol{\theta})
    =&
    \frac{1}{2}
    \sum_{j\in\mathcal{S}}
    w_j
    \left[
      \frac{
      \ln y_j^{\mathrm{sim}}(\boldsymbol{\theta})
      -
      \ln y_j^{\mathrm{exp}}
      }{
      \sigma_j
      }
    \right]^2
    \\
    &+
    w_W
    \sum_{k\in\mathcal{D}}
    W_1\!\left[
      p_k^{\mathrm{sim}}(R;\boldsymbol{\theta}),
      p_k^{\mathrm{exp}}(R)
    \right]
    +
    \Phi_{\mathrm{phys}} .
\end{split}
\label{eq:dt_neg_log_likelihood}
\end{equation}
Here $\mathcal{S}$ is the set of scalar observables, $\mathcal{D}$ is the set
of distributional observables, and $w_j$ and $w_W$ are user-assigned or
hierarchically inferred weights.

The physical penalty is
\begin{equation}
    \Phi_{\mathrm{phys}}
    =
    \lambda_{\mathrm{FP}}\delta_{\mathrm{FP}}^2
    +
    \lambda_{\mathrm{gas}}\sum_\mu \delta_{g_\mu}^2
    +
    \lambda_{\mathrm{bounds}}\Phi_{\mathrm{bounds}} .
    \label{eq:dt_physical_penalty}
\end{equation}
The quantities $\delta_{\mathrm{FP}}$ and $\delta_{g_\mu}$ are conservation
diagnostics for signed Frenkel-pair content and gas inventory.  In accepted
high-fidelity calculations these errors should remain below $10^{-6}$; values
above $10^{-3}$ should trigger rejection of the run unless the loss is due to a
modeled external sink.

%===================================================================
\subsection{Multi-Fidelity Simulation and Emulation}
\label{app:dt_multifidelity}
%===================================================================

Direct Bayesian calibration on the full RadCluster simulator is not practical
for $p\simeq16$--20 uncertain inputs.  The twin therefore uses a two-level
multi-fidelity hierarchy.  The low-fidelity simulator,
$\mathcal{M}_{\mathrm{LF}}$, uses coarse size-bin moments, a reduced cluster grid,
and a shorter dose horizon, but preserves verified solver tolerances.  The
high-fidelity simulator, $\mathcal{M}_{\mathrm{HF}}$, uses the full bin resolution,
the full loop-conversion and network model, and is integrated to reactor or
fusion-relevant dose.

A space-filling design is generated in the prior parameter domain,
\begin{equation}
    \left\{
      \boldsymbol{\theta}^{(k)}
    \right\}_{k=1}^{N_{\mathrm{LF}}}
    \sim
    \mathrm{Sobol}(p),
    \qquad
    N_{\mathrm{LF}}\simeq 20p\text{--}50p .
\end{equation}
A smaller high-fidelity design is then selected,
\begin{equation}
    N_{\mathrm{HF}}\simeq 3p\text{--}8p,
\end{equation}
using either a nested space-filling subset or active-learning criteria.

The low- and high-fidelity outputs are fused using an autoregressive
multi-fidelity Gaussian-process emulator,
\begin{equation}
    \widehat{\mathcal{M}}_{\mathrm{HF}}(\boldsymbol{\theta},\boldsymbol{\xi})
    =
    \rho
    \widehat{\mathcal{M}}_{\mathrm{LF}}(\boldsymbol{\theta},\boldsymbol{\xi})
    +
    \delta(\boldsymbol{\theta},\boldsymbol{\xi}),
    \label{eq:dt_multifidelity_gp}
\end{equation}
where $\delta$ is a Gaussian-process correction.  The emulator reports both
mean predictions and predictive variances, so emulator uncertainty enters the
likelihood through Eq.~\eqref{eq:dt_uncertainty_sum}.

%===================================================================
\subsection{Sensitivity Screening and Identifiability}
\label{app:dt_sensitivity}
%===================================================================

Before full calibration, the Sobol ensemble is used to compute first-order,
interaction, and total-effect sensitivity indices,
\begin{equation}
    S_i
    =
    \frac{\mathrm{Var}_{\theta_i}
    \left[\mathbb{E}(Y\mid\theta_i)\right]}{\mathrm{Var}(Y)},
    \qquad
    S_i^T
    =
    1-
    \frac{\mathrm{Var}_{\boldsymbol{\theta}_{\sim i}}
    \left[\mathbb{E}(Y\mid\boldsymbol{\theta}_{\sim i})\right]}
    {\mathrm{Var}(Y)} .
\end{equation}
Parameters with small total-effect indices over all measured observables are
not eliminated from the physics model, but they are demoted during inverse
calibration.  This prevents weakly identifiable parameters from absorbing
experimental scatter or model discrepancy.  The active calibration subset is
therefore
\begin{equation}
    \boldsymbol{\theta}_{\mathrm{act}}
    =
    \left\{
      \theta_i:
      \max_Y S_i^T(Y) > S_{\mathrm{tol}}
    \right\},
\end{equation}
while inactive parameters are either fixed at nominal values or marginalized
with their prior uncertainty.

%===================================================================
\subsection{Sequential Updating and Digital-Twin Assimilation}
\label{app:dt_sequential_update}
%===================================================================

When measurements arrive sequentially along a dose trajectory, the twin is
updated recursively.  Let
\begin{equation}
    \mathbf{z}_k
    =
    H(\mathbf{x}_k)
    +
    \boldsymbol{\epsilon}_k
\end{equation}
be a measurement at dose step $k$.  The prior at step $k$ is the posterior
from step $k-1$ propagated through RadCluster:
\begin{equation}
    p(\boldsymbol{\theta},\mathbf{x}_k\mid\mathbf{z}_{1:k-1})
    =
    \int
    p(\mathbf{x}_k\mid\mathbf{x}_{k-1},\boldsymbol{\theta})
    p(\boldsymbol{\theta},\mathbf{x}_{k-1}\mid\mathbf{z}_{1:k-1})
    d\mathbf{x}_{k-1} .
\end{equation}
The update is
\begin{equation}
    p(\boldsymbol{\theta},\mathbf{x}_k\mid\mathbf{z}_{1:k})
    \propto
    p(\mathbf{z}_k\mid\mathbf{x}_k,\boldsymbol{\theta})
    p(\boldsymbol{\theta},\mathbf{x}_k\mid\mathbf{z}_{1:k-1}) .
\end{equation}
For high-dimensional ensembles, this update can be approximated by ensemble
Kalman inversion or an ensemble smoother.  This converts RadCluster from a
one-shot calibrated code into a living digital twin whose parameter posterior
and microstructural state are continually revised as new irradiation data are
obtained.

%===================================================================
\subsection{Active Refinement and Experiment Selection}
\label{app:dt_active_refinement}
%===================================================================

The surrogate should be refined only where additional RadCluster runs or new
experiments most reduce uncertainty in the predicted microstructure.  Define
an information-per-cost acquisition function
\begin{equation}
    \mathcal{A}(\boldsymbol{\theta},\boldsymbol{\xi})
    =
    \frac{
      \mathcal{I}(\boldsymbol{\theta},\boldsymbol{\xi})
    }{
      C_{\mathrm{CPU}}(\boldsymbol{\theta},\boldsymbol{\xi})
    },
\end{equation}
where $\mathcal{I}$ may be the expected reduction in posterior entropy, the
expected reduction in predictive variance at fusion conditions, or the expected
improvement in validation likelihood.  The next simulation is chosen by
\begin{equation}
    (\boldsymbol{\theta}_{\mathrm{next}},\boldsymbol{\xi}_{\mathrm{next}})
    =
    \arg\max_{\boldsymbol{\theta},\boldsymbol{\xi}}
    \mathcal{A}(\boldsymbol{\theta},\boldsymbol{\xi}) .
\end{equation}
The same criterion can be used to recommend new irradiation experiments:
select the temperature, dose, dose rate, and gas-production ratio that most
reduce the uncertainty in predictions at DEMO-relevant conditions.

%===================================================================
\subsection{Workflow Summary}
\label{app:dt_workflow_summary}
%===================================================================

Figure~\ref{fig:dt_workflow} summarizes the offline construction of the
RadCluster digital twin.  Figure~\ref{fig:dt_assimilation} shows the online or
sequential assimilation mode.

\begin{figure}[htbp]
\centering
\begin{tikzpicture}[
    node distance=1.8cm and 1.7cm,
    block/.style={draw, rounded corners, align=center, minimum width=3.0cm,
                  minimum height=0.95cm, font=\small},
    smallblock/.style={draw, rounded corners, align=center, minimum width=2.6cm,
                  minimum height=0.85cm, font=\small},
    arrow/.style={->, thick}
]
\node[block] (priors) {Physics-based\\priors and bounds\\$p(\boldsymbol{\theta})$};
\node[block, right=of priors] (design) {Sobol/LHS design\\$\{\boldsymbol{\theta}^{(k)}\}$};
\node[block, right=of design] (runs) {RadCluster ensemble\\LF and HF runs};
\node[block, below=of runs] (features) {Observation operator\\$H[\mathbf{x}]$\\extracts $N_\ell,\bar d_\ell,N_c,\bar d_c,p(R)$};
\node[block, left=of features] (emu) {Multi-fidelity GP\\emulator\\$\widehat{\mathcal{M}}_{\rm HF}$};
\node[block, left=of emu] (cal) {Bayesian calibration\\$p(\boldsymbol{\theta}|\mathcal{D}_{\rm exp})$};
\node[smallblock, below=of cal] (targets) {Appendix G targets\\trend bands + histograms};
\node[smallblock, below=of emu] (cons) {Conservation checks\\$\delta_{\rm FP},\delta_g$};
\node[block, right=of features] (pred) {Posterior prediction\\fusion-relevant\\microstructure};
\node[block, above=of pred] (active) {Active refinement\\information/cost};

\draw[arrow] (priors) -- (design);
\draw[arrow] (design) -- (runs);
\draw[arrow] (runs) -- (features);
\draw[arrow] (features) -- (emu);
\draw[arrow] (emu) -- (cal);
\draw[arrow] (targets) -- (cal);
\draw[arrow] (cons) -- (cal);
\draw[arrow] (cal) -- (pred);
\draw[arrow] (pred) -- (active);
\draw[arrow] (active) |- (runs);
\end{tikzpicture}
\caption{Offline construction of the RadCluster digital twin.  Physical priors define the
uncertain parameter space, a space-filling design generates RadCluster ensembles, the
observation operator extracts experimental observables, a multi-fidelity emulator enables
Bayesian calibration, and active refinement requests additional simulations where they
most reduce predictive uncertainty per unit computational cost.}
\label{fig:dt_workflow}
\end{figure}

\begin{figure}[htbp]
\centering
\begin{tikzpicture}[
    node distance=1.65cm and 1.6cm,
    block/.style={draw, rounded corners, align=center, minimum width=3.0cm,
                  minimum height=0.9cm, font=\small},
    arrow/.style={->, thick}
]
\node[block] (post0) {Prior/posterior at dose $d_{k-1}$\\$p(\boldsymbol{\theta},\mathbf{x}_{k-1}|\mathbf{z}_{1:k-1})$};
\node[block, right=of post0] (prop) {RadCluster propagation\\$\mathbf{x}_{k-1}\rightarrow\mathbf{x}_k$};
\node[block, right=of prop] (pred) {Predict measurement\\$\hat{\mathbf{z}}_k=H(\mathbf{x}_k)$};
\node[block, below=of pred] (data) {New experiment\\$\mathbf{z}_k$};
\node[block, left=of data] (update) {Bayesian/EnKF update\\assimilate residual};
\node[block, left=of update] (post1) {Updated twin\\$p(\boldsymbol{\theta},\mathbf{x}_k|\mathbf{z}_{1:k})$};
\node[block, below=of update] (decision) {Prediction, validation,\\or new experiment selection};

\draw[arrow] (post0) -- (prop);
\draw[arrow] (prop) -- (pred);
\draw[arrow] (pred) -- (data);
\draw[arrow] (data) -- (update);
\draw[arrow] (update) -- (post1);
\draw[arrow] (post1) |- (decision);
\draw[arrow] (decision) -| (prop);
\end{tikzpicture}
\caption{Sequential digital-twin assimilation.  New measurements at dose step $k$ are
compared with RadCluster predictions through the observation operator.  The residual is
used to update the joint posterior over model parameters and microstructural state, after
which the twin is propagated to the next dose or used to design the next informative
experiment.}
\label{fig:dt_assimilation}
\end{figure}

%===================================================================
\subsection{Acceptance Criteria for a Validated Twin}
\label{app:dt_acceptance_criteria}
%===================================================================

A RadCluster digital twin should be accepted only if it satisfies the following
conditions.
\begin{enumerate}
    \item \textbf{Statistical consistency:} scalar predictions fall within the
    appropriate Appendix~G log-normal prediction bands, preferably within
    $\pm 2s$ in log space.
    \item \textbf{Distributional consistency:} predicted loop and cavity size
    distributions match measured histograms within an agreed Wasserstein
    tolerance.
    \item \textbf{Conservation consistency:} high-fidelity accepted runs satisfy
    $\delta_{\mathrm{FP}}<10^{-6}$ and $\delta_{g_\mu}<10^{-6}$, unless a
    modeled external sink explicitly removes material.
    \item \textbf{Predictive validation:} posterior predictive intervals cover
    withheld validation data at approximately their nominal credibility level.
    \item \textbf{Extrapolation discipline:} fusion-relevant predictions are
    reported as posterior predictive distributions, not deterministic curves.
\end{enumerate}

Once these criteria are met, the calibrated twin can be used to predict
microstructure under fusion-relevant conditions,
\begin{equation}
    T=300\text{--}550^\circ\mathrm{C},
    \qquad
    d=0.1\text{--}150~\mathrm{dpa},
    \qquad
    \eta_{\mathrm{He}}=10\text{--}15~\mathrm{appm/dpa},
    \qquad
    \eta_{\mathrm{H}}=40\text{--}50~\mathrm{appm/dpa},
\end{equation}
with credible intervals for loop density, loop size, cavity density, cavity
size, swelling, loop-character fractions, and the probability of entering a
bimodal cavity-growth regime.

%===================================================================
\subsection{Implementation Checklist}
\label{app:dt_implementation_checklist}
%===================================================================

A practical implementation should contain the following files or modules:
\begin{enumerate}
    \item \texttt{parameters.yaml}: symbols, units, nominal values, bounds,
    prior types, and active/inactive flags for every uncertain parameter.
    \item \texttt{experiments.yaml} or a database reader for
    \texttt{RadiationDatabase.xlsx}: irradiation conditions, observables,
    uncertainty estimates, grouping variables, and distributional data.
    \item \texttt{run\_ensemble.py}: low- and high-fidelity RadCluster ensemble
    execution with conservation diagnostics and solver-status logging.
    \item \texttt{extract\_observables.py}: maps cluster distributions to
    $N_\ell$, $\bar d_\ell$, $N_c$, $\bar d_c$, swelling, loop-character
    fractions, and histograms.
    \item \texttt{surrogate.py}: multi-fidelity Gaussian-process training,
    prediction, and emulator uncertainty estimation.
    \item \texttt{calibrate.py}: posterior inference using MCMC, sequential
    Monte Carlo, variational inference, or ensemble Kalman inversion.
    \item \texttt{active\_refinement.py}: selection of additional RadCluster
    runs or proposed experiments by information gain per computational cost.
    \item \texttt{predict\_fusion.py}: posterior predictive simulation under
    fusion-relevant temperature, dose, dose-rate, He/dpa, and H/dpa histories.
\end{enumerate}

%===================================================================
\subsection{BibTeX Entries for Digital-Twin/UQ Citations}
\label{app:dt_bibtex_entries}
%===================================================================

The following entries may be added to the project bibliography if they are not
already present.

\begin{verbatim}
@article{Thelen2022DigitalTwinUQ,
  author  = {Thelen, Adam and Zhang, Xiaoge and Fink, Olga and Lu, Yan and
             Ghosh, Sayan and Youn, Byeng D. and Todd, Michael D. and
             Mahadevan, Sankaran and Hu, Chao and Hu, Zhen},
  title   = {A comprehensive review of digital twin---part 2: roles of
             uncertainty quantification and optimization, a battery digital
             twin, and perspectives},
  journal = {Structural and Multidisciplinary Optimization},
  year    = {2022},
  volume  = {66},
  pages   = {1},
  doi     = {10.1007/s00158-022-03410-x}
}

@article{Wang2020AMUQ,
  author  = {Wang, Zhuo and Jiang, Chen and Liu, Pengwei and Yang, Wenhua and
             Chen, Lei},
  title   = {Uncertainty quantification and reduction in metal additive
             manufacturing},
  journal = {npj Computational Materials},
  year    = {2020},
  volume  = {6},
  pages   = {175},
  doi     = {10.1038/s41524-020-00444-x}
}

@article{Kim2022BayesianAM,
  author  = {Kim, Jee Yun and Garcia, David and Zhu, Yunhui and Higdon, David M.
             and others},
  title   = {A Bayesian learning framework for fast prediction and uncertainty
             quantification of additively manufactured multi-material components},
  journal = {Journal of Materials Processing Technology},
  year    = {2022},
  volume  = {303},
  pages   = {117528},
  doi     = {10.1016/j.jmatprotec.2022.117528}
}

@article{Robertson2026VDMN,
  author  = {Robertson, Andreas E. and Inman, Samuel B. and Lenau, Ashley T. and
             Lebensohn, Ricardo A. and Shin, Dongil and Boyce, Brad L. and
             Dingreville, Remi M.},
  title   = {Microstructure-based variational deep material networks for robust
             uncertainty quantification in materials digital twins},
  journal = {npj Computational Materials},
  year    = {2026},
  doi     = {10.1038/s41524-026-02054-5}
}
\end{verbatim}
